%% ma: L10-B14-0088
%% lop: 10
%% chuong: 5
%% bai: 14
%% dang: 10.14.12
%% loai: DS
%% muc: [NB, TH, TH, VD]
%% dap_an: "ĐSĐĐ"
%% nguon: "chuyen-de-cac-so-dac-trung-cua-mau-so-lieu-khong-ghep-nhom-toan-10-chuong-trinh-moi.pdf (GV. Phan Nhật Linh) (P05), Bài 03 Dạng 2 Phần II Câu 5"
%% kien_thuc: "Lớp 10 Bài 13 (số trung bình, trung vị, tứ phân vị, mốt), Bài 14 (khoảng biến thiên, khoảng tứ phân vị, phương sai, độ lệch chuẩn, giá trị bất thường)"
%% trang_thai: chinh-thuc
%% ghi_chu: "Bảng gốc có thêm số 157 (11 số cho 10 tỉnh) mà lời giải gốc không dùng, đã bỏ; ý d) gốc s ≈ 2,57 (do làm tròn phương sai thành 6,6), sửa thành 2,56. Tự duyệt theo yêu cầu giáo viên 22/10/2026."
\begin{ex}
Tỉ lệ trẻ suy dinh dưỡng (\%, tính theo cân nặng ứng với độ tuổi) của $10$ tỉnh thuộc Đồng bằng sông Hồng được cho như sau (theo Tổng cục Thống kê):
\begin{center}
\begin{tabular}{|c|c|c|c|c|c|c|c|c|c|}\hline
$5{,}5$ & $13{,}8$ & $10{,}2$ & $12{,}2$ & $11{,}0$ & $7{,}4$ & $11{,}4$ & $13{,}1$ & $12{,}5$ & $13{,}4$\\\hline
\end{tabular}
\end{center}
Xét tính đúng sai của các khẳng định sau:
\choiceTF{\True Khoảng biến thiên của mẫu số liệu là $R=8{,}3$}{Trung vị của mẫu số liệu là $Q_2=9{,}2$}{\True Số trung bình của mẫu số liệu là $\overline x=11{,}05$}{\True Độ lệch chuẩn của mẫu số liệu (làm tròn đến hàng phần trăm) là $s\approx2{,}56$}
\loigiai{Sắp xếp: $5{,}5$; $7{,}4$; $10{,}2$; $11{,}0$; $11{,}4$; $12{,}2$; $12{,}5$; $13{,}1$; $13{,}4$; $13{,}8$.

a) Đúng: $R=13{,}8-5{,}5=8{,}3$.

b) Sai: $Q_2=\dfrac{11{,}4+12{,}2}2=11{,}8$.

c) Đúng: $\overline x=\dfrac{110{,}5}{10}=11{,}05$.

d) Đúng: $s^2=6{,}5685$; $s\approx2{,}56$.}
\end{ex}
